\documentclass[10pt,a4paper]{amsart}
\usepackage[margin=19mm]{geometry}
\usepackage{amsmath,amssymb,mathtools}
\usepackage[scr=rsfs,cal=euler]{mathalfa}
\usepackage{xfrac}
\usepackage[unicode,hidelinks,bookmarks=false]{hyperref}
\newcommand{\ie}{i.e.\ }
\newcommand{\cf}{cf.\ }
\newcommand{\resp}{respectively\ }
\newcommand{\Subsection}{\subsection}
\providecommand{\nobreakdash}{\nobreak\hbox{-}}
\setlength{\emergencystretch}{2em}
\multlinegap0pt
\usepackage{bm}
\usepackage{bbm}
\usepackage{dsfont}
\usepackage{xspace}
\usepackage{cancel}
\usepackage{mathtools}
\usepackage[shortlabels]{enumitem}
\usepackage{amsthm}
\makeatletter
\@ifundefined{claim}{\newtheorem{claim}{Claim}}{}
\@ifundefined{thm}{\newtheorem{thm}{Theorem}}{}
\makeatother
\newcommand{\Z}{\mathbb{Z}}
\newcommand{\rk}{{\rm rk}\,}
\title[Special alternating links of minimal unlinking number]{Special alternating links of minimal unlinking number}
\author{Duncan McCoy, JungHwan Park}
\date{Source: arXiv:2603.10732v2 (CC BY 4.0)}
\begin{document}
\maketitle

\noindent\emph{Source and licence.} Special alternating links of minimal unlinking number. Version arXiv:2603.10732v2 (CC BY 4.0), distributed under the Creative Commons Attribution 4.0 International licence, as recorded in the arXiv licence metadata for this exact version and shown on the abstract page https://arxiv.org/abs/2603.10732v2. Full rights notice: licenses/arxiv-2603.10732-CC-BY-4.0.txt.\par\smallskip

\noindent\textit{Editorial scope.} Counted unit: the Claim whose source label is \texttt{claim1} in the article (source0.tex lines 355--358, source label \texttt{claim1}), delivered together with the article's own complete original proof of it (source0.tex lines 359--379, one \texttt{proof} environment of 21 lines). No mathematics is written, repaired or re-derived: statement and proof are the article's own text line for line. Required context retained uncounted: thm:lattice\_obstruction (lines 264--271); it:all\_coords (lines 267--270). Every definition, display and label of those lines is retained unchanged. Omitted from this package and not redistributed: 1--263, 271--354, 380--564. Nothing inside a retained excerpt is omitted or altered: every retained body is delivered exactly as the article's lines are written, so no omission receipt of any kind accompanies this package. Citations: the retained excerpts carry no citation command at all, so no bibliography entry of the article is transcribed and no reference is redistributed. Reference adaptation: none was needed and none was made; every reference inside the delivered unit resolves inside it, so no unresolved reference or citation is produced. Printed slips kept literal: no printed word, symbol, hypothesis or number of the counted statement or of its proof was altered. Layout and class: the article's own class is \texttt{amsart} with packages including babel, amsmath, amsthm, amsfonts, enumerate, graphicx, color, xy, thmtools, thm-restate, overpic, amssymb, tikz, quiver; the wrapper is the lane's pinned \texttt{amsart} editorial wrapper at 10pt on A4, which restates the theorem-like environments the retained text needs and adds the article's own macro definitions that the retained text needs (\texttt{\textbackslash Z}, \texttt{\textbackslash rk}), with extra wrapper packages (bm, bbm, dsfont, xspace, cancel, mathtools, enumitem). The delivered numbering is the unit's own. Assets: no retained span contains a figure environment, a graphics-inclusion command, a PDF-page inclusion, a table or a TikZ picture; no image file is copied into this package, the package declares no asset dependency and no image file is redistributed with it. Licence: version arXiv:2603.10732v2 of this article is distributed under the Creative Commons Attribution 4.0 International licence, as recorded in the arXiv licence metadata for this exact version and shown on the abstract page https://arxiv.org/abs/2603.10732v2. A full rights notice is delivered at licenses/arxiv-2603.10732-CC-BY-4.0.txt. Characters: every retained excerpt is pure ASCII, so the delivered engine, XeLaTeX with its default Latin Modern text font, reports no missing character.
\begin{thm}\label{thm:lattice_obstruction}
Let $L$ be a non-split $k$-component alternating link with $\sigma(L)\leq 0$. Let $D$ be a reduced alternating diagram for $L$ with positive-definite Goeritz lattice $\Lambda_D$. Let 
\[p=\frac{-\sigma(L)+k-1}{2} \quad\text{and}\quad n=\rk \Lambda_D-\sigma(L).\] If $c_4(L)=p$, then there is a lattice embedding of $\Lambda_D$ into $\Z^n$ such that the following conditions are satisfied:
\begin{enumerate}[(i)]
    \item\label{it:all_coords} for all unit vectors $e\in \Z^n$, there exists $v$ in the image of $\Lambda_D$ such that $v\cdot e\neq 0$ and
    \item\label{it:u_disjoint_supports} there exists an orthonormal basis $e_1,\dots, e_n$ for $\Z^n$ such that for all $i=1, \dots, p$, every $v$ in the image of $\Lambda_D$ satisfies $v\cdot e_{i}=v\cdot e_{i+p}$.
\end{enumerate}
\end{thm}

\begin{claim}\label{claim1}
    For any basis vector $e_i\in \Z^n$, there are two regions $v_\alpha, v_\beta$ such that 
    \[v_\alpha \cdot e_i = +1,\quad v_\beta \cdot e_i=-1\] and for all other $j\neq \alpha,\beta$ we have $v_j\cdot e_i=0$.
\end{claim}

\begin{proof}[Proof of Claim]
By Theorem~\ref{thm:lattice_obstruction}\eqref{it:all_coords}, the lattice $\Lambda_D$ uses every coordinate, and therefore
\begin{equation}\label{eq:1}
\sum_{j=0}^r \bigl|v_j\cdot e_i\bigr|>0
\end{equation}
for all $i=1,\dots,n$. On the other hand, the fact that $\sum_{j=0}^r v_j=0$ implies that the left-hand side of \eqref{eq:1} is even. Thus,
\begin{equation}\label{eq:2}
\sum_{j=0}^r \bigl|v_j\cdot e_i\bigr|\ge 2
\end{equation}
for all $i=1,\dots,n$. The number of crossings of $D$ can be computed via
\begin{equation}\label{eq:3}
2n=\sum_{j=0}^r v_j\cdot v_j
=\sum_{j=0}^r\sum_{i=1}^n \bigl|v_j\cdot e_i\bigr|^2
\ge \sum_{i=1}^n\sum_{j=0}^r \bigl|v_j\cdot e_i\bigr|.
\end{equation}
If we sum the inequalities in \eqref{eq:2} over all $i$ and compare with \eqref{eq:3}, then we obtain
\begin{equation}\label{eq:4}
\sum_{j=0}^r \bigl|v_j\cdot e_i\bigr|=2
\end{equation}
for all $i=1,\dots,n$. Invoking once again that $\sum_{j=0}^r v_j=0$, the claim follows from \eqref{eq:4}.
\end{proof}

\end{document}
